\documentclass[12pt, a4paper]{article}

\usepackage{amsmath}
\usepackage{amssymb}
\usepackage{titlesec}
\usepackage{tcolorbox}
\usepackage{enumitem}
\usepackage[bottom=3cm]{geometry}
\usepackage[hidelinks]{hyperref}
\tcbuselibrary{breakable}
\tcbuselibrary{skins}

\newtcolorbox{prob}[1]{colback=gray!5!white, colframe=gray!75!black,
    title=\textbf{Exercise #1}}

\newtcolorbox{sol}{
    breakable,
    colback=white,
    colframe=white,
    frame hidden,
    left=3mm, right=3mm,
    boxrule=0mm,
    top=0mm, bottom=0mm,
    parbox=false,
    before upper={\textbf{Solution:}\par\medskip}
}

\begin{document}

\begin{prob}{27.1}
A photon counter is pointed at a remote star for one minute, in order to infer the rate of photons arriving at the counter per minute, $\lambda$. Assuming the number of photons collected $r$ has a Poisson distribution with mean $\lambda$,
\[
P(r \mid \lambda) = \exp(-\lambda)\frac{\lambda^r}{r!},
\]
and assuming the improper prior $P(\lambda) = 1/\lambda$, make Laplace approximations to the posterior distribution
\begin{enumerate}[label=(\alph*)]
    \item over $\lambda$
    \item over $\log \lambda$. [Note the improper prior transforms to $P(\log \lambda) = \text{constant}$.]
\end{enumerate}
\end{prob}
\begin{sol}
\end{sol}
\bigskip

\begin{prob}{27.2}
Use Laplace's method to approximate the integral
\[
Z(u_1, u_2) = \int_{-\infty}^{\infty} \mathrm{d}a \, f(a)^{u_1} \left(1 - f(a)\right)^{u_2},
\]
where $f(a) = 1/(1 + e^{-a})$ and $u_1, u_2$ are positive. Check the accuracy of the approximation against the exact answer (23.29, p.\,316) for $(u_1, u_2) = (\tfrac{1}{2}, \tfrac{1}{2})$ and $(u_1, u_2) = (1, 1)$. Measure the error $(\log Z_P - \log Z_Q)$ in bits.
\end{prob}
\begin{sol}
\end{sol}
\bigskip

\begin{prob}{27.3}
\emph{Linear regression.} $N$ datapoints $\{(x^{(n)}, t^{(n)})\}$ are generated by the experimenter choosing each $x^{(n)}$, then the world delivering a noisy version of the linear function
\[
y(x) = w_0 + w_1 x,
\]
\[
t^{(n)} \sim \text{Normal}\left(y(x^{(n)}), \sigma_\nu^2\right).
\]
Assuming Gaussian priors on $w_0$ and $w_1$, make the Laplace approximation to the posterior distribution of $w_0$ and $w_1$ (which is exact, in fact) and obtain the predictive distribution for the next datapoint $t^{(N+1)}$, given $x^{(N+1)}$. (See MacKay (1992a) for further reading.)
\end{prob}
\begin{sol}
\end{sol}
\bigskip

\end{document}
